\chapter{Import external images}

\eyebrow{importing photos while preserving pairing}

\vspace{4pt}
For photos already acquired (taken with another camera, in another session), the
\textbf{Import images} object brings them in without breaking pairing.

\section{Two ways to import}

A folder is selected; the app scans it, computes a \tele{sha256} for each image
(so that duplicates are detected), and reads EXIF. Two import modes are then
available:

\begin{enumerate}[leftmargin=1.4em,itemsep=4pt]
  \item \textbf{Import all as new carcasses.} Each photo becomes a new carcass,
        with the physical tag taken from the filename. The real data is entered
        later in the carcass Overview. Repeated filenames receive a numeric suffix
        automatically, so tags remain unique.
  \item \textbf{Pair to an existing carcass.} Per item, a photo is reconciled to a
        carcass already in the batch. This applies when several photos belong to
        the same animal.
\end{enumerate}

\begin{note}[title=Pairing is preserved either way]
In both modes, the carcass is created or chosen at the moment of import, so the
image is not orphaned. Duplicates already in the dataset are skipped by hash.
\end{note}

This converts the earlier failure mode, orphan images later given invented links,
into an explicit step resolved by the operator at import time.
